\chapter{Linear Regression, from a Line to the Normal Equations}
\companion{StatQuest: Linear Regression, Clearly Explained}{\ytid{7ArmBVF2dCs}}

Linear regression is the first model interviewers drill and the one they come back to in Round 4
(``what are its assumptions? what if homoscedasticity fails?''). It is worth knowing completely: every idea
in it (a loss, a gradient, a closed form, assumptions, diagnostics, regularisation) reappears in every other
model.

\section{The model in words and symbols}

\subsection{One feature}
We believe the target $y$ is roughly a straight-line function of one feature $x$:
\[ \hat y = w x + b. \]
$w$ is the \textbf{slope} (how much $\hat y$ rises when $x$ rises by one unit) and $b$ the \textbf{intercept}
(the prediction when $x = 0$).

\begin{example}[: reading the numbers]
A fitted salary model: $\widehat{\text{salary}} = 1.8\times\text{experience} + 3.5$ (lakh). A fresher ($x=0$) is
predicted 3.5 lakh. Each extra year adds 1.8 lakh. Someone with 4 years: $1.8\cdot4+3.5 = 10.7$ lakh.
\end{example}

\subsection{Many features}
With $d$ features $x_1,\dots,x_d$ (experience, a city indicator, number of skills, \dots):
\[ \hat y = w_1x_1 + w_2x_2 + \dots + w_dx_d + b = \mathbf{w}^\top\mathbf{x} + b. \]
The symbol $\mathbf{w}^\top\mathbf{x}$ is a \textbf{dot product}: multiply matching entries and add them up.
For $\mathbf{w} = (2, -1)$ and $\mathbf{x} = (3, 4)$, $\mathbf{w}^\top\mathbf{x} = 2\cdot3 + (-1)\cdot4 = 2$.

Stack $n$ examples as rows of a matrix $X$ ($n$ rows, $d$ columns, plus a column of 1s to carry the
intercept). Then all predictions at once are $\hat{\mathbf y} = X\mathbf w$.

\subsection{Why it is called ``linear''}
``Linear'' means linear in the \emph{parameters} $\mathbf w$, not in $x$. The model
$\hat y = w_1 x + w_2 x^2 + w_3\log x + b$ is still linear regression: we simply created new features
$x^2$ and $\log x$. This is how linear regression fits curves.

\section{Choosing the line: least squares}

\subsection{Residuals and the loss}
For each example, the \textbf{residual} is the error $e_i = y_i - \hat y_i$. We want residuals small. We
cannot just add them up (positive and negative errors would cancel), so we square them:
\[ \text{SSE} = \sum_{i=1}^n (y_i - \hat y_i)^2, \qquad \text{MSE} = \frac{1}{n}\text{SSE}. \]
\textbf{Ordinary least squares (OLS)} chooses $w, b$ to minimise this.

Why squares and not absolute values?
\begin{enumerate}
\item Squares are smooth, so calculus gives a clean closed-form answer.
\item Squares punish large errors heavily (an error of 4 costs 16, four errors of 1 cost 4). Often
  appropriate, sometimes a weakness (outliers dominate).
\item If the noise is Gaussian, least squares is exactly maximum likelihood (Chapter 12). That is the deep
  reason.
\end{enumerate}

\subsection{Solving it with one feature, step by step}
A \textbf{derivative} tells you the slope of a function: where a smooth function is at its lowest point, its
slope is zero. So we take the derivative of SSE with respect to $b$ and $w$ and set each to zero.
\begin{itemize}
\item Setting $\partial\text{SSE}/\partial b = -2\sum(y_i - wx_i - b) = 0$ gives $b = \bar y - w\bar x$.
  ($\bar x$ means the average of the $x$ values.) So the line always passes through $(\bar x, \bar y)$.
\item Substituting and setting $\partial\text{SSE}/\partial w = 0$ gives
  \[ w = \frac{\sum (x_i - \bar x)(y_i - \bar y)}{\sum (x_i - \bar x)^2} = \frac{\Cov(x,y)}{\Var(x)} = r\,\frac{s_y}{s_x}, \]
  where $r$ is the correlation and $s_x, s_y$ the standard deviations.
\end{itemize}

\begin{example}[: fit a line to four Noida flats]
Area $x$ (hundreds of sq ft): 5, 7, 9, 11. Price $y$ (lakh): 30, 40, 44, 54.
$\bar x = 8$, $\bar y = 42$. Deviations: $x - \bar x = -3,-1,1,3$; $y-\bar y = -12,-2,2,12$.
$\sum(x-\bar x)(y-\bar y) = 36+2+2+36 = 76$. $\sum(x-\bar x)^2 = 9+1+1+9 = 20$.
$w = 76/20 = 3.8$ lakh per 100 sq ft. $b = 42 - 3.8\cdot8 = 11.6$.
Line: $\hat y = 3.8x + 11.6$. A 1{,}000 sq ft flat ($x=10$): $49.6$ lakh.
Predictions $30.6, 38.2, 45.8, 53.4$; residuals $-0.6, 1.8, -1.8, 0.6$; they sum to $0$, as they must.
\end{example}

\subsection{Many features: the normal equations}
In matrix form, SSE $= \norm{\mathbf y - X\mathbf w}^2$. Its gradient (the vector of all partial derivatives)
is $-2X^\top(\mathbf y - X\mathbf w)$. Setting it to zero:
\[ X^\top X\,\mathbf w = X^\top\mathbf y \quad\Longrightarrow\quad \mathbf w = (X^\top X)^{-1}X^\top\mathbf y. \]
These are the \textbf{normal equations}. They require $X^\top X$ to be invertible, which fails when one
feature is an exact combination of others (perfect multicollinearity) or when there are more features than rows.

\subsubsection{The geometric picture (a favourite interview follow-up)}
Think of $\mathbf y$ as one point in $n$-dimensional space (one coordinate per example). Every possible
prediction $X\mathbf w$ lies in the flat subspace spanned by the columns of $X$ (the \textbf{column space}).
Least squares picks the point in that subspace closest to $\mathbf y$: the \textbf{orthogonal projection}. So
the residual vector $\mathbf y - X\hat{\mathbf w}$ is perpendicular to every column of $X$. That is exactly what
$X^\top(\mathbf y - X\mathbf w) = 0$ says. With a column of ones in $X$, perpendicular to it means the residuals
sum to zero.

\subsubsection{Cost and alternatives}
Forming $X^\top X$ costs about $nd^2$ operations and inverting it about $d^3$. Fine for $d$ in the thousands.
For huge $d$ or streaming data, use gradient descent (Chapter 5). In practice libraries use a QR or SVD
decomposition rather than an explicit inverse, which is numerically safer.

\section{Measuring the fit}

\subsection{$R^2$: fraction of variance explained}
Compare our model with the dumbest possible model, ``always predict $\bar y$''.
\[ R^2 = 1 - \frac{\text{SS}_{res}}{\text{SS}_{tot}} = 1 - \frac{\sum(y_i - \hat y_i)^2}{\sum(y_i - \bar y)^2}. \]
$R^2 = 1$: perfect; $R^2 = 0$: no better than the mean; on test data $R^2$ can be negative (worse than the mean).
For simple regression with an intercept, $R^2 = r^2$.

\begin{example}[: $R^2$ for the flats]
SS$_{res} = 0.36 + 3.24 + 3.24 + 0.36 = 7.2$. SS$_{tot} = 144+4+4+144 = 296$. $R^2 = 1 - 7.2/296 = 0.976$.
\end{example}

\subsection{Adjusted $R^2$}
Training $R^2$ never decreases when you add a feature (OLS can always give it weight 0). Adjusted $R^2$
penalises features: $\bar R^2 = 1 - (1-R^2)\frac{n-1}{n-d-1}$. It goes down if a new feature adds less than
chance would.

\subsection{Error metrics}
\begin{itemize}
\item \textbf{MAE} $=\frac1n\sum|y-\hat y|$: in target units, robust to outliers; minimised by the median.
\item \textbf{RMSE} $=\sqrt{\text{MSE}}$: in target units, punishes large errors; minimised by the mean.
\item \textbf{MAPE} $=\frac1n\sum|y-\hat y|/|y|$: percentage, scale-free, explodes near $y=0$ and penalises
  over-prediction more than under-prediction.
\end{itemize}

\section{Assumptions: what OLS needs, and for what}

A crucial distinction interviewers love: some assumptions are needed for the \emph{coefficients to be
right on average} (unbiased), others only for \emph{p-values and confidence intervals to be valid}.

\begin{center}\small
\begin{tabularx}{\linewidth}{l X X X}
\toprule
Assumption & Plain meaning & How to check & If violated\\
\midrule
\textbf{L}inearity & the average of $y$ moves linearly with the features (after any transforms) & residuals vs
fitted plot shows a curve & add $x^2$, $\log x$, interactions, or use a non-linear model\\
\textbf{I}ndependence & errors of different rows are unrelated & Durbin--Watson; residuals vs time & time-series
models, clustered standard errors, mixed models\\
\textbf{N}ormality & errors are bell-shaped & Q--Q plot, Shapiro--Wilk & matters only for small-sample inference;
transform $y$, bootstrap, or robust methods\\
\textbf{E}qual variance (homoscedasticity) & spread of errors is the same everywhere & residual plot shows a
funnel; Breusch--Pagan test & transform $y$ (log), weighted least squares, robust SEs, or a GLM\\
No perfect multicollinearity & no feature is an exact combination of others & VIF, correlation matrix & drop or
combine features, PCA, Ridge\\
Exogeneity & errors are uncorrelated with features (no omitted confounder) & domain reasoning & coefficients
biased; add the confounder, instrumental variables\\
\bottomrule
\end{tabularx}
\end{center}

Mnemonic: \textbf{LINE} plus ``no multicollinearity''. For unbiased coefficients you need linearity,
exogeneity and no perfect multicollinearity. Independence and homoscedasticity make OLS the \emph{best}
(minimum-variance) linear unbiased estimator (the Gauss--Markov theorem) and make the usual standard errors
correct. Normality is only for exact small-sample t- and F-tests.

\subsection{Heteroscedasticity, slowly}
``Hetero'' means different, ``scedastic'' means spread. Salaries are the textbook case: freshers' salaries vary by
$\pm1$ lakh, senior leaders' by $\pm20$ lakh. A residual plot fans out like a funnel.
\begin{itemize}
\item \textbf{What breaks:} coefficients stay unbiased, but OLS standard errors (and hence p-values) are wrong,
  and OLS gives too much influence to the noisiest points.
\item \textbf{Fixes:}
  \begin{enumerate}
  \item \emph{Transform $y$}: model $\log(\text{salary})$. Multiplicative noise ($\pm10\%$) becomes additive,
    roughly constant noise. (This is what the reported interviewer wanted.)
  \item \emph{Weighted least squares}: give each point weight $1/\sigma_i^2$, so noisy points count less.
  \item \emph{Robust (``sandwich'', HC) standard errors}: keep OLS coefficients, fix the inference.
  \item \emph{A GLM} with a variance that grows with the mean, such as Gamma regression with a log link
    (what the candidate proposed; Chapter 4 explains it fully).
  \end{enumerate}
\end{itemize}

\subsection{Multicollinearity and VIF}
If two features move together (``years of experience'' and ``age''), the model cannot tell which one deserves
credit. Predictions stay fine, but individual coefficients become unstable: they can swing wildly or flip sign
between samples. The \textbf{variance inflation factor} of feature $j$ comes from regressing $x_j$ on all other
features and taking its $R_j^2$:
\[ \text{VIF}_j = \frac{1}{1 - R_j^2}. \]
VIF $=1$: no collinearity; VIF $>5$ or $>10$: worrying. The standard error of $w_j$ is inflated by $\sqrt{\text{VIF}_j}$.

\section{Regularised linear regression}
\begin{itemize}
\item \textbf{Ridge}: minimise SSE $+\lambda\norm{\mathbf w}^2$. Closed form $\mathbf w = (X^\top X + \lambda I)^{-1}X^\top\mathbf y$.
  Adding $\lambda$ to the diagonal makes the matrix always invertible, which is why Ridge cures multicollinearity.
\item \textbf{Lasso}: minimise SSE $+\lambda\sum|w_j|$. No closed form; solved by coordinate descent. Sets some
  weights exactly to zero.
\item \textbf{Elastic net}: both penalties.
\end{itemize}
The intercept is not penalised, and features must be standardised first (Chapter 1).

\section{Outliers, leverage and influence}
\begin{itemize}
\item An \textbf{outlier} has an unusual $y$ given its $x$ (large residual).
\item A \textbf{high-leverage point} has an unusual $x$ (far from the other $x$ values). It has the power to
  tilt the line.
\item An \textbf{influential point} actually does tilt it: high leverage \emph{and} a large residual.
  \textbf{Cook's distance} measures how much all fitted values change when that point is deleted.
\end{itemize}
Remedies: check whether it is a data error; use robust losses (Huber, MAE); or keep it if it is real and
important. Deleting genuine extreme values is a mistake the InfoEdge coding round specifically tested.

\begin{practical}[: where linear regression lives at InfoEdge]
\begin{itemize}
\item \textbf{Salary engine} (Data Factory slide): estimating a fair salary range for a title, city and
  experience. Log-salary linear models with many one-hot features are a strong, interpretable baseline.
\item \textbf{99acres price estimates}: price per square foot by locality, floor, age, amenities.
\item As a \textbf{diagnostic tool}: coefficients answer ``how much does one more year of experience add, holding
  city fixed?'', which product teams ask constantly.
\end{itemize}
\end{practical}

\stuck{StatQuest: R-squared, Clearly Explained}{\ytid{2AQKmw14mHM}}
\section{Interview questions}

\begin{qa}{\pyq{INT, Round 4 follow-up} What are the assumptions of linear regression?}
\oneline{Linearity in the parameters, independent errors, constant error variance (homoscedasticity), normally
distributed errors, no perfect multicollinearity, and errors uncorrelated with the features (exogeneity).}
\why
I like to group them by what they buy us.
\begin{enumerate}
\item \textbf{For coefficients to be correct on average (unbiased):}
  \begin{itemize}
  \item \emph{Linearity:} $\E[y\mid x]$ is a linear combination of the features we included. Note: features can
    be transformed ($x^2$, $\log x$), so this is about the specification.
  \item \emph{Exogeneity:} the error has mean zero given the features. Violated when an omitted variable affects
    both $x$ and $y$ (omit ``company tier'' from a salary model and the ``IIT degree'' coefficient absorbs part of
    its effect).
  \item \emph{No perfect multicollinearity:} otherwise the coefficients are not even uniquely defined.
  \end{itemize}
\item \textbf{For OLS to be the most efficient linear unbiased estimator, and for standard errors to be right:}
  \emph{independence} of errors and \emph{homoscedasticity} (Gauss--Markov).
\item \textbf{For exact small-sample t-tests and confidence intervals:} \emph{normality} of errors. With large $n$
  the central limit theorem makes this matter little.
\end{enumerate}
\worked For each I state a check: residual vs fitted plot (linearity, homoscedasticity), Durbin--Watson or
residuals vs time (independence), Q--Q plot (normality), VIF (multicollinearity).
\pushes \emph{``Do the features need to be normally distributed?''} No; the assumptions are about the errors.
Binary features are perfectly fine. \emph{``Does $y$ need to be normal?''} Not marginally; only the errors given
$x$.
\pitfall Saying ``the data must be normal''. Also listing assumptions without saying which consequences follow
from violating each.
\end{qa}

\begin{qa}{\pyq{INT, Round 4} Homoscedasticity is violated in your model. What do you do?}
\oneline{First confirm it with a residual plot or Breusch--Pagan test; then either transform the target (log or
Box--Cox), use weighted least squares, keep OLS with heteroscedasticity-robust standard errors, or switch to a GLM
whose variance grows with the mean.}
\why
\textbf{What goes wrong.} The coefficients are still unbiased, so predictions of the mean are fine. But OLS treats
every point as equally reliable, so it is no longer the most efficient estimator, and the usual standard-error
formula (which assumes one common $\sigma^2$) is wrong: p-values and confidence intervals cannot be trusted.
Prediction intervals are also wrong: too wide for small values, too narrow for large ones.

\textbf{Fix 1, transform $y$.} If the noise is proportional to the level (salaries vary by about $\pm15\%$), take
logs: $\log(y) = \log(\text{level}) + \log(1\pm0.15)$, and the second term has roughly constant spread. Box--Cox
searches over power transforms $y^\lambda$. Back-transforming needs care: $\exp(\E[\log y]) \ne \E[y]$; multiply by
a smearing factor such as $\exp(\hat\sigma^2/2)$.

\textbf{Fix 2, weighted least squares.} If the variance of point $i$ is $\sigma_i^2$, minimise
$\sum (y_i - \hat y_i)^2/\sigma_i^2$. Noisy points get less weight. Usually $\sigma_i$ is modelled (for instance
proportional to $\hat y_i$).

\textbf{Fix 3, robust standard errors.} Keep OLS, but compute standard errors with the White/HC formula that does
not assume constant variance. Good when the goal is inference and the coefficients themselves are fine.

\textbf{Fix 4, a GLM.} Model $y$ with a distribution whose variance is tied to the mean: Gamma (variance $\propto$
mean$^2$, constant coefficient of variation) or Poisson for counts (variance $=$ mean), usually with a log link.
This models the mean on the original scale, avoiding the back-transformation problem.
\worked Salary vs experience: residual spread is about 1 lakh at 1 year and 8 lakh at 15 years. Fit
$\log(\text{salary})$; the residual plot becomes a flat band; the coefficient $0.08$ means ``each year adds about 8\%''.
\pushes \emph{``Which would you choose?''} For prediction with a multiplicative structure, log or Gamma-log GLM; for
explaining effects with valid p-values, robust SEs.
\pitfall Claiming heteroscedasticity biases the coefficients. It does not (unless it comes from misspecification).
\end{qa}

\begin{qa}{\pyq{INT: multicollinearity and VIF} What is multicollinearity, why is it a problem, and how do you
detect and fix it?}
\oneline{It is strong correlation among features; it makes individual coefficients unstable and their standard
errors large, though predictions are largely unaffected; detect it with VIF and fix it by dropping, combining,
PCA or Ridge.}
\why
If $x_1$ and $x_2$ almost always move together, many combinations of $(w_1, w_2)$ give nearly the same predictions
($w_1 + w_2 \approx$ constant). The data cannot pin down the split, so small changes in the sample swing the
coefficients, sometimes flipping signs. Mathematically, $X^\top X$ is nearly singular, so its inverse (which sets
the coefficient variances) has huge entries.

\textbf{Detect:} pairwise correlations miss three-way relations, so use VIF$_j = 1/(1-R_j^2)$. Also the condition
number of $X$. \textbf{Fix:} drop one of the pair; combine them (an index); use PCA; use Ridge (which adds $\lambda I$ and
stabilises the inverse); or collect data where they vary independently.
\worked Regressing ``years of experience'' on ``age'' and ``graduation year'' gives $R^2 = 0.96$, so VIF $= 25$: the
standard error of the experience coefficient is $5\times$ what it would be without collinearity.
\pushes \emph{``Does it matter for a random forest?''} Prediction is fine, but importance gets split between the
correlated features, so each looks less important than it is.
\end{qa}

\begin{qa}{Derive the OLS slope and intercept for a single feature.}
\oneline{Set the derivatives of the squared-error sum with respect to $b$ and $w$ to zero; you get $b = \bar y - w\bar x$
and $w = \sum(x-\bar x)(y-\bar y)/\sum(x-\bar x)^2$.}
\why
$S(w,b) = \sum(y_i - wx_i - b)^2$.
\[ \frac{\partial S}{\partial b} = -2\sum(y_i - wx_i - b) = 0 \;\Rightarrow\; \sum y_i = w\sum x_i + nb \;\Rightarrow\; b = \bar y - w\bar x. \]
Substitute $b$: residual $= (y_i - \bar y) - w(x_i - \bar x)$.
\[ \frac{\partial S}{\partial w} = -2\sum x_i\big[(y_i-\bar y) - w(x_i-\bar x)\big] = 0. \]
Because $\sum \bar x[(y_i-\bar y) - w(x_i - \bar x)] = 0$, we can replace $x_i$ by $(x_i - \bar x)$:
\[ \sum(x_i-\bar x)(y_i-\bar y) = w\sum(x_i-\bar x)^2 \;\Rightarrow\; w = \frac{S_{xy}}{S_{xx}}. \]
The second derivative $2\sum(x_i - \bar x)^2 > 0$, so this is a minimum.
\worked See the Noida flats example: $w = 76/20 = 3.8$, $b = 11.6$.
\end{qa}

\begin{qa}{Derive the normal equations and give their geometric meaning.}
\oneline{Minimising $\norm{\mathbf y - X\mathbf w}^2$ gives $X^\top X\mathbf w = X^\top\mathbf y$; geometrically, the fitted
vector is the orthogonal projection of $\mathbf y$ onto the column space of $X$, so residuals are perpendicular to every
feature column.}
\why Expand $L = (\mathbf y - X\mathbf w)^\top(\mathbf y - X\mathbf w) = \mathbf y^\top\mathbf y - 2\mathbf w^\top X^\top\mathbf y + \mathbf w^\top X^\top X\mathbf w$.
Gradient: $-2X^\top\mathbf y + 2X^\top X\mathbf w = 0$. If $X^\top X$ is invertible, $\mathbf w = (X^\top X)^{-1}X^\top\mathbf y$.
The \textbf{hat matrix} $H = X(X^\top X)^{-1}X^\top$ maps $\mathbf y$ to $\hat{\mathbf y}$; it is a projection ($H^2 = H$),
and its diagonal entries are the leverages. Its trace equals the number of parameters.
\pushes \emph{``What if $X^\top X$ is not invertible?''} Infinitely many solutions fit equally well; use the
pseudo-inverse (minimum-norm solution, via SVD), drop redundant columns, or add a Ridge penalty.
\end{qa}

\begin{qa}{What does $R^2$ mean? Can it be negative? Why do we need adjusted $R^2$?}
\oneline{$R^2$ is the fraction of the target's variance the model explains relative to predicting the mean; on new
data it can be negative; adjusted $R^2$ corrects for the fact that training $R^2$ never falls when features are added.}
\why $R^2 = 1 - \text{SS}_{res}/\text{SS}_{tot}$. On training data with an intercept, OLS can always do at least as
well as the mean, so $0 \le R^2 \le 1$. On test data a bad model can have SS$_{res} > $ SS$_{tot}$. Adjusted
$R^2 = 1 - (1-R^2)\frac{n-1}{n-d-1}$ rises only if a new feature improves fit more than expected by chance.
\worked $n = 50$, $d = 4$, $R^2 = 0.80$: adjusted $= 1 - 0.2\cdot49/45 = 0.782$.
\pitfall ``$R^2 = 0.9$ means the model is good/causal.'' A high $R^2$ can come from leakage or a trend, and a low
$R^2$ can still be very useful (noisy human behaviour).
\end{qa}

\begin{qa}{\deep Can linear regression fit a curve? Can it fit $y = a e^{bx}$?}
\oneline{Yes: ``linear'' means linear in the parameters, so adding $x^2$, $x^3$ or $\log x$ features fits curves; an
exponential with multiplicative noise becomes linear after taking logs.}
\why A polynomial model $\hat y = w_0 + w_1x + w_2x^2$ is a weighted sum of fixed functions of $x$, so OLS solves it
with the same normal equations. For $y = ae^{bx}\varepsilon$ (multiplicative noise $\varepsilon$), $\log y = \log a + bx +
\log\varepsilon$: regress $\log y$ on $x$; the slope estimates $b$ and the intercept estimates $\log a$. If the
noise is \emph{additive} ($y = ae^{bx} + \varepsilon$), log-transforming distorts it; then use non-linear least
squares or a GLM with a log link (Poisson/Gamma), which models $\log\E[y] = \log a + bx$ directly.
\worked Number of applications to a job vs days since posting decays exponentially: regressing $\log(\text{applies}+1)$
on days gives slope $-0.12$, i.e. about $11\%$ fewer applies per extra day ($e^{-0.12} = 0.887$).
\end{qa}

\begin{qa}{\deep Linear regression on a 0/1 target: what goes wrong? Is it ever acceptable?}
\oneline{Predictions can fall outside $[0,1]$, errors are heteroscedastic by construction, and far-away easy points
tilt the line; it is acceptable (the ``linear probability model'') for quick effect estimates near the middle of the
range, but logistic regression is the right default.}
\why For a binary $y$, the error variance is $p(1-p)$, which depends on $x$, so homoscedasticity fails automatically.
Squared loss also penalises a correctly classified but extreme point (predicted 1.7 for a true 1) and pulls the line
toward it, moving the 0.5 threshold. Economists still use the linear probability model because its coefficients read
directly as ``change in probability'' and with robust SEs it is fine near $p\approx0.5$.
\end{qa}

\begin{qa}{What is the difference between correlation and the regression slope? When are they equal?}
\oneline{Slope $= r\cdot s_y/s_x$; it carries units and direction of prediction, correlation is unit-free and symmetric;
they are equal when both variables are standardised.}
\why Regressing $y$ on $x$ gives slope $r s_y/s_x$; regressing $x$ on $y$ gives $r s_x/s_y$; the two lines differ
unless $|r| = 1$. Their product is $r^2$. After converting both variables to z-scores, the slope is exactly $r$.
\worked If $r = 0.6$, $s_y = 5$ lakh, $s_x = 2$ years, the slope is $0.6\cdot5/2 = 1.5$ lakh per year.
\end{qa}

\begin{qa}{Explain Ridge regression and why it handles multicollinearity.}
\oneline{Ridge adds $\lambda\norm{\mathbf w}^2$ to the squared error; its solution $(X^\top X + \lambda I)^{-1}X^\top\mathbf y$ always
exists, and the added $\lambda$ stabilises directions in which the data carry little information.}
\why With collinear features $X^\top X$ has a tiny eigenvalue along the direction where features cancel ($w_1 -
w_2$). Inverting divides by that tiny number, so coefficients explode along that direction. Ridge adds $\lambda$ to
every eigenvalue, so the smallest becomes at least $\lambda$, and the explosion stops. Equivalently it shrinks the
coefficients most along low-variance principal directions. The price is a little bias.
\worked One standardised feature, $\sum x_i^2 = 10$, $\sum x_iy_i = 20$: OLS $w = 2$; Ridge $\lambda = 10$: $w =
20/(10+10) = 1$.
\end{qa}

\begin{qa}{Lasso vs Ridge vs Elastic net: when would you use each?}
\oneline{Lasso when you believe few features matter and want selection; Ridge when many features each matter a little
or are correlated; Elastic net when you want sparsity but have correlated groups.}
\why Lasso's L1 penalty zeroes weights (feature selection) but with correlated features it keeps one arbitrarily and
can be unstable, and it can select at most $n$ features when $d > n$. Ridge never zeroes but is stable. Elastic net
($\alpha$L1 $+ (1-\alpha)$L2) groups correlated features and still sparsifies.
\inpractice For a salary model with 5{,}000 one-hot skill features, Elastic net gives a compact, stable model whose
non-zero skills are explainable to product teams (``Kubernetes adds 6\%'').
\end{qa}

\begin{qa}{What are leverage, outliers and influential points? How do you detect them?}
\oneline{Leverage is unusualness in $x$, an outlier is unusualness in $y$ given $x$, and an influential point is one whose
removal changes the fit a lot; leverage comes from the hat matrix diagonal, influence from Cook's distance.}
\why A high-leverage point that lies on the trend does no harm and even improves precision. A high-leverage point
off the trend acts like a long lever and pulls the line. Cook's distance combines residual size and leverage;
values above $4/n$ are commonly flagged. Studentised residuals flag outliers.
\pushes \emph{``Would you delete it?''} Only if it is an error (a price of Rs.~1, a negative area). If it is a real
luxury villa, keep it, and consider a robust loss or a separate model for that segment.
\end{qa}

\begin{qa}{\deep You fit a regression of salary on ``number of certifications'' and get a negative coefficient.
Does that mean certifications reduce salary?}
\oneline{Not necessarily: it is a correlation conditional on the other included features, and it can be driven by
confounding (juniors collect more certificates) or by collinearity with other features.}
\why A coefficient measures the association of a feature with the target \emph{holding the other features fixed}.
If experience is not in the model, ``many certificates'' may simply indicate a junior person (omitted-variable
bias). If experience is in the model and highly collinear, the sign can flip between samples. Causal claims need
either an experiment or a careful causal design (controls for confounders, instrumental variables, difference in
differences).
\worked Without experience: coefficient $-0.4$ lakh/cert. Adding experience: $+0.3$ lakh/cert. This is Simpson's
paradox in regression form.
\end{qa}

\begin{qa}{Gradient descent or normal equations for linear regression: which and when?}
\oneline{Normal equations (via QR/SVD) for small to medium $d$ because they are exact and need no tuning; gradient
descent (or SGD) for very large $d$, very large or streaming $n$, or when you add penalties without closed forms.}
\why Normal equations cost about $O(nd^2 + d^3)$ and need the whole data in memory. Gradient descent costs
$O(nd)$ per pass, can stream mini-batches, but needs a learning rate, feature scaling and many iterations.
With $d = 100$ and $n = 10^6$: normal equations are instant. With $d = 10^6$ sparse text features: SGD.
\end{qa}

\begin{qa}{\deep Why do we minimise squared error and not absolute error? When would absolute error be better?}
\oneline{Squared error is smooth, has a closed form, and is the maximum-likelihood loss for Gaussian noise;
absolute error is better when there are outliers or heavy-tailed noise, and it targets the median instead of the mean.}
\why The value $c$ that minimises $\sum(y_i - c)^2$ is the mean; the one that minimises $\sum|y_i - c|$ is the
median. So OLS estimates the conditional mean and least-absolute-deviation regression estimates the conditional
median. With one data-entry error of Rs.~10 crore in a property dataset, the mean (and OLS) moves a lot; the median
barely moves. Huber loss is a compromise: squared for small errors, absolute for large ones.
\worked Values $\{3, 4, 5, 6, 100\}$: mean $= 23.6$, median $= 5$.
\end{qa}

\begin{qa}{Interpret the coefficient $0.08$ in a regression of $\log(\text{salary})$ on years of experience.
What if both sides were logged?}
\oneline{Each extra year multiplies salary by $e^{0.08}\approx1.083$, about $+8.3\%$; in a log--log model the coefficient
is an elasticity: a 1\% increase in $x$ gives about a $\beta$\% increase in $y$.}
\why $\log y = \beta_0 + 0.08x$; raising $x$ by 1 adds 0.08 to $\log y$, i.e. multiplies $y$ by $e^{0.08}$. For small
$\beta$, $e^\beta - 1 \approx \beta$. In $\log y = \beta_0 + \beta\log x$, $\frac{d\log y}{d\log x} = \beta$ (elasticity).
\worked Price per sq ft vs area in log--log with $\beta = -0.15$: a 10\% larger flat has about 1.5\% lower price per
sq ft (bulk discount).
\end{qa}

\begin{qa}{How would you check the assumptions of a fitted linear model in Python?}
\oneline{Plot residuals against fitted values and against time, draw a Q--Q plot, compute VIFs, and run
Breusch--Pagan and Durbin--Watson tests using statsmodels.}
\why
\begin{lstlisting}
import statsmodels.api as sm
from statsmodels.stats.diagnostic import het_breuschpagan
from statsmodels.stats.outliers_influence import variance_inflation_factor
Xc = sm.add_constant(X)
m = sm.OLS(y, Xc).fit()
print(m.summary())                     # coefficients, p-values, R2, Durbin-Watson
resid, fitted = m.resid, m.fittedvalues
# plt.scatter(fitted, resid)  -> look for curves (linearity) and funnels (heteroscedasticity)
sm.qqplot(resid, line="s")            # normality
print(het_breuschpagan(resid, Xc))    # small p-value -> heteroscedastic
vifs = [variance_inflation_factor(Xc.values, i) for i in range(1, Xc.shape[1])]
m_robust = sm.OLS(y, Xc).fit(cov_type="HC3")   # robust standard errors
\end{lstlisting}
\end{qa}

\begin{qa}{\deep Your linear model's training $R^2$ is 0.35. A colleague says linear regression is useless here.
How do you respond?}
\oneline{A low $R^2$ only says much of the variance is unexplained; the model can still be useful if its coefficients are
reliable or its ranking or average predictions are good enough for the decision, and the gap may be closable with better
features rather than a different model.}
\why Human behaviour (who applies, who buys) is noisy; an $R^2$ of 0.35 can be excellent. Ask: (1) what is the
business metric? For ranking leads, ordering matters, not exact values; (2) is it underfitting? Compare with a
gradient-boosted model: if it also gets 0.37, the ceiling is the noise, not linearity; if it gets 0.60, add
non-linear terms or switch; (3) are features missing? Domain features (price-band concentration in 99acres) often
beat model changes.
\end{qa}

\begin{qa}{What does the intercept mean, and should you ever remove it?}
\oneline{The intercept is the prediction when all features are zero; remove it only when theory demands the line pass
through the origin, because forcing it biases the slope otherwise.}
\why Without an intercept, residuals no longer need to sum to zero and $R^2$ can be misleading (some software reports
an uncentred $R^2$). If features are centred, the intercept equals $\bar y$, which is easy to interpret. A case for
no intercept: cost $=$ price $\times$ quantity, where zero quantity truly means zero cost.
\end{qa}

\begin{qa}{What is the Gauss--Markov theorem, in plain words?}
\oneline{Among all estimators that are linear in $\mathbf y$ and unbiased, OLS has the smallest variance, provided errors have
mean zero, constant variance and are uncorrelated; it is ``BLUE''.}
\why ``Best linear unbiased estimator.'' It does not require normality. It does not say OLS beats biased estimators:
Ridge is biased but can have lower total error, which is the bias--variance trade-off again.
\end{qa}

\begin{qa}{\deep Can you do linear regression when you have more features than rows ($d > n$)?}
\oneline{Not with plain OLS, because $X^\top X$ is singular and infinitely many weight vectors fit the data exactly; with
Ridge, Lasso, PCA regression or the minimum-norm solution you can.}
\why With 1{,}000 skill features and 200 salaries, any 200-dimensional pattern can be fitted perfectly, so the
training error is zero and meaningless. Ridge makes the problem well-posed; Lasso selects at most $n$ features;
the pseudo-inverse picks the smallest-norm exact fit (which, interestingly, is what gradient descent from zero
converges to). Always validate, since all of these can still overfit.
\end{qa}
